\begin{remark}[McQuillan's variant]
\label{remark-mcquillan}
There is a variant of the construction of formal schemes due to
McQuillan, see \cite{McQuillan}.
He suggests a slight weakening of the condition of admissibility.
Namely, recall that an admissible topological ring is a complete
(and separated by our conventions) topological ring $A$
which is linearly topologized such that there exists an
ideal of definition: an
open ideal $I$ such that any neighbourhood of $0$ contains $I^n$
for some $n \geq 1$.
McQuillan works with what we will call {\it weakly admissible}
topological rings. A weakly admissible topological ring $A$ is a
complete (and separated by our conventions) topological ring
which is linearly topologized such that there exists an
{\it weak ideal of definition}: an open ideal $I$ such that
for all $f \in I$ we have
$f^n \to 0$ for $n \to \infty$. Similarly to the admissible case,
if $I$ is a weak ideal of definition and $J \subset A$ is an
open ideal, then $I \cap J$ is a weak ideal of definition.
Thus the weak ideals of definition form a fundamental system of
open neighbourhoods of $0$ and
one can proceed along much the same route as above
to define a larger category of formal schemes based
on this notion. The analogues of Lemmas \ref{lemma-fully-faithful} and
\ref{lemma-formal-scheme-sheaf-fppf}
still hold in this setting (with the same proof).
\end{remark}